% U8 WS3 -- structure, pKa, salts, common ion. Blocks 5-6.
\documentclass{apchem-worksheet}

\apcunit{8}
\apcunittitle{Acids and Bases}
\apcdoctitle{Worksheet 3 \textbullet\ Structure, p$K_a$, and Salts}
\apcsubtitle{CED 8.6--8.7, 8.4 \textbullet\ trace every salt back to its parent acid and base}

\begin{document}
\wsheader[p$K_a = -\log K_a$; lower p$K_a$ means a stronger acid
\textbullet\ $K_b = K_w/K_a$ \textbullet\ $K_a$: \ce{HF}
$7.2\times10^{-4}$, \ce{HC2H3O2} $1.8\times10^{-5}$, \ce{HCN}
$6.2\times10^{-10}$; $K_b(\ce{NH3}) = 1.8\times10^{-5}$.]

\problem[]{\ced{8.7} p$K_a$ practice:}
\begin{parts}
  \item p$K_a$ of acetic acid: \answerblank[2.4cm]{4.74}
  \item p$K_a$ of \ce{HCN}: \answerblank[2.4cm]{9.21}
  \item p$K_a$ of \ce{HF}: \answerblank[2.4cm]{3.14}
  \item Which is the strongest acid of the three, and how does p$K_a$ show
        it? \answerlines[2]{\ce{HF} --- the lowest p$K_a$ corresponds to the
        largest $K_a$ and therefore the greatest extent of ionization.}
\end{parts}

\problem[]{\ced{8.6} Structure and acid strength.}
\begin{parts}
  \item Rank \ce{HF}, \ce{HCl}, \ce{HBr}, \ce{HI} by increasing strength
        and give the controlling factor.
        \answerlines[2]{\ce{HF} $<$ \ce{HCl} $<$ \ce{HBr} $<$ \ce{HI}. The
        H--X bond weakens down the group, so the proton is released more
        readily; bond strength outweighs electronegativity.}
  \item Explain why \ce{HF} is weakest despite fluorine being the most
        electronegative element. \answerlines[2]{The H--F bond is by far the
        strongest of the four, and holding the proton tightly matters more
        than pulling electron density away from it.}
  \item Rank \ce{HClO}, \ce{HClO2}, \ce{HClO3}, \ce{HClO4} and give both
        reasons. \answerlines[3]{Strength increases with the number of
        oxygens. Each additional electronegative oxygen withdraws electron
        density from the O--H bond, weakening it, and spreads the negative
        charge of the conjugate base, stabilizing it.}
  \item Rank \ce{HOI}, \ce{HOBr}, \ce{HOCl}.
        \answerblank[5cm]{\ce{HOI} $<$ \ce{HOBr} $<$ \ce{HOCl}}
\end{parts}

\problem[]{\ced{8.7} Classify each salt solution and name the responsible
ion:}
\begin{parts}
  \item \ce{NaCl}: \answerblank[7.1cm]{neutral --- neither ion hydrolyses}
  \item \ce{NH4Cl}: \answerblank[6.5cm]{acidic --- \ce{NH4+}}
  \item \ce{NaCN}: \answerblank[6.5cm]{basic --- \ce{CN-}}
  \item \ce{KNO3}: \answerblank[6.5cm]{neutral}
  \item \ce{KF}: \answerblank[6.5cm]{basic --- \ce{F-}}
\end{parts}

\problem[]{\ced{8.7} Explain, in conjugate-pair terms, why \ce{NaCl}
solutions are neutral but \ce{NaCN} solutions are basic.
\answerlines[3]{Both contain \ce{Na+}, from the strong base \ce{NaOH},
which does not react. \ce{Cl-} comes from the strong acid \ce{HCl}, so it is
a negligible base. \ce{CN-} comes from the very weak acid \ce{HCN}, so it
readily accepts a proton from water and generates \ce{OH-}.}}

\problem[]{\ced{8.7} Calculate the pH of \SI{0.20}{\Molar}
\ce{NaC2H3O2}.}
\begin{parts}
  \item Write the hydrolysis equation.
        \answerlines[1]{\ce{C2H3O2^-(aq) + H2O(l) <=> HC2H3O2(aq) +
        OH^-(aq)}}
  \item Calculate $K_b$. \answerblank[3.2cm]{$5.6\times10^{-10}$}
  \item Calculate the pH. \workspace[2.4cm]
        \keyonly{\begin{apcanswer}$x = \sqrt{(5.6\times10^{-10})(0.20)} =
        1.06\times10^{-5} = [\ce{OH-}]$; pOH $= 4.98$;
        pH $= \mathbf{9.02}$\end{apcanswer}}
\end{parts}

\problem[]{\ced{8.7} Calculate the pH of \SI{0.30}{\Molar} \ce{NaF}.
\workspace[2.8cm]
\keyonly{\begin{apcanswer}$K_b = 1.0\times10^{-14}/7.2\times10^{-4} =
1.4\times10^{-11}$; $x = \sqrt{(1.4\times10^{-11})(0.30)} =
2.0\times10^{-6}$; pOH $= 5.69$; pH $=
\mathbf{8.31}$\end{apcanswer}}}

\problem[]{\ced{8.4} The common-ion effect. For \SI{1.0}{\Molar} \ce{HF}
alone, $[\ce{H+}] = 2.7\times10^{-2}$ and the acid is 2.7\% ionized.}
\begin{parts}
  \item Calculate $[\ce{H+}]$ in a solution \SI{1.0}{\Molar} in \ce{HF} and
        \SI{1.0}{\Molar} in \ce{NaF}. \workspace[2cm]
        \keyonly{\begin{apcanswer}$[\ce{H+}] = K_a[\ce{HF}]/[\ce{F-}] =
        (7.2\times10^{-4})(1.0/1.0) = 7.2\times10^{-4}$;
        pH $= \mathbf{3.14}$\end{apcanswer}}
  \item Percent ionization now: \answerblank[2.6cm]{0.072\%}
  \item Explain the suppression using Le Ch\^atelier.
        \answerlines[3]{\ce{F-} is a product of the \ce{HF} ionization
        equilibrium. Adding it shifts the equilibrium to the left, so less
        \ce{HF} ionizes and $[\ce{H+}]$ falls --- here by a factor of about
        37.}
\end{parts}

\problem[]{\ced{8.4} Buffer or not? Give a one-phrase reason.}
\begin{parts}
  \item \ce{HC2H3O2} $+$ \ce{NaC2H3O2}
        \answerblank[8.9cm]{buffer --- weak acid and its conjugate base}
  \item \ce{HCl} $+$ \ce{NaCl}
        \answerblank[7cm]{not --- \ce{HCl} is strong, \ce{Cl-} inert}
  \item \ce{NH3} $+$ \ce{NH4NO3}
        \answerblank[8.9cm]{buffer --- weak base and its conjugate acid}
  \item \ce{NaOH} $+$ \ce{NaC2H3O2}
        \answerblank[8cm]{not --- no weak conjugate pair present}
\end{parts}

\begin{spiralstrip}{Unit 8 \textbullet\ blocks 1--4}
  \item For a weak base, ICE $x$ is: \answerblank[2.2cm]{$[\ce{OH-}]$}
  \item pH of \SI{0.100}{\Molar} \ce{NH3}: \answerblank[2cm]{11.13}
  \item Which ionization dominates for \ce{H3PO4}?
        \answerblank[2.4cm]{the first}
  \item pH of \SI{0.10}{\Molar} acetic acid: \answerblank[2cm]{2.87}
  \item \ce{NH4Cl} solutions are: \answerblank[2cm]{acidic}
\end{spiralstrip}

\end{document}
